%%%%%%%%%%  scan idx 339  |  folio 326  %%%%%%%%%%
\textbf{3.3.}
Soient \(K\), \(K'\) deux \(\lambda\)-anneaux, \(f:K\to K'\) un homomorphisme d'anneaux. D'après le définition 3.2.ii), \(f\) est un \(\lambda\)-homomorphisme si et seulement si pour tout \(i\in\mathbb N\), et tout \(x\in K\), \(f(\gamma^i(x))=\gamma^i(f(x))\), ou encore \(\gamma_t(f(x))=\Gg(f)\circ\gamma_t(x)\).

\textbf{3.4.} De la relation 3.2.ii), on déduit les relations :
\begin{equation}\tag{3.4.1}
  \gamma^0(x)=1,
  \qquad \gamma^1(x)=x,
  \qquad \gamma^n(x+y)=\sum_{p+q=n}\gamma^p(x)\gamma^q(y).
\end{equation}
Elles montrent que l'homomorphisme \(\gamma_t\) définit sur \(K\) une nouvelle structure de pré-\(\lambda\)-anneau ; nous allons voir que ce n'est pas en général une structure de \(\lambda\)-anneau. Si on désigne par \(\rho_t\) la substitution sur l'anneau des séries formelles en \(t\) à coefficients dans \(K\) définie par \(\rho_t(t)=t/1-t\), on a :
\[
  \gamma_t=\rho_t\circ\lambda_t.
\]
L'automorphisme \(\rho_t\) est fonctoriel par rapport à \(K\), et définit un automorphisme du foncteur \(\Ghat\) considéré comme foncteur en groupes abéliens ; à l'aide de cet automorphisme, on peut, par transport de structure, définir sur \(\Ghat\) une nouvelle structure de foncteur en \(\lambda\)-anneaux, qui sera telle que pour tout \(\lambda\)-anneau \(K\), l'homomorphisme \(\gamma_t:K\to\Ghat(K)\) soit un \(\lambda\)-homomorphisme ; nous allons expliciter la multiplication et les opérations \(\lambda^i\) de cette deuxième structure.

\textbf{3.5.} La multiplication est définie par l'homomorphisme de foncteurs :
\[
  \Ghat\times\Ghat\xrightarrow{(\rho^{-1},\rho^{-1})}
  \Ghat\times\Ghat\xrightarrow{\circ}\Ghat\xrightarrow{\rho}\Ghat ,
\]
où le signe \(\circ\) note la multiplication définie en 2.3. Cette nouvelle multiplication sera notée \(v\). Elle est comme d'habitude définie par sa valeur sur \((1+Xt,1+Yt)\)~; on a donc :
